\title{Direct Preference Optimization: Your Language Model is Secretly a Reward Model}

\begin{abstract}
Large-scale unsupervised language models (LMs) can acquire extensive world knowledge and some capacity for reasoning, yet precisely shaping their behavior remains challenging because of the lack of supervision during training. To address this, existing strategies typically gather human feedback in the form of comparative judgments on model outputs and adjust the LM to better reflect these preferences—commonly using reinforcement learning from human feedback (RLHF). However, the RLHF pipeline introduces significant complexity and instability, as it first requires training a reward model to represent human preferences, followed by reinforcement learning to optimize the unsupervised LM according to this reward signal while preserving the model’s original characteristics. In this work, we propose a novel formulation for the reward model within RLHF, allowing the optimal policy corresponding to human preferences to be derived analytically in closed form. This insight leads to a straightforward approach where the conventional RLHF problem can be addressed by optimizing a simple classification objective. The proposed method, termed \textit{Direct Preference Optimization} (DPO), is both robust and efficient, dispensing with the need for language model sampling or extensive hyperparameter search during fine-tuning. Experimental evaluations demonstrate that DPO aligns LMs with human preferences at least as well as, or often better than, existing approaches. Importantly, DPO surpasses RLHF approaches based on PPO in managing generation sentiment and achieves equal or superior performance on summarization and single-turn dialogue benchmarks, all while being considerably easier to use and train.
\end{abstract}

\section{Introduction}
Large-scale unsupervised language models (LMs), trained on immense corpora, often develop unexpected proficiencies~\citep{chowdhery2022palm, brown2020language, touvron2023llama,bubeck2023sparks}. Nevertheless, because these models are built upon human-generated content reflecting a broad spectrum of aims, values, and levels of expertise, not all aspects present in the training data are equally desirable to replicate. For instance, in the context of AI coding assistants, it is beneficial for the model to \textit{recognize} prevalent coding errors to facilitate their correction. Yet, during code generation, preference should be given to the exceptional coding practices, albeit rare, encapsulated within the dataset. Analogously, while awareness of a widely held misconception—embraced by 50\% of people—may be useful for the model, it is undesirable for the model to echo this incorrect belief in half of its responses to relevant questions. Thus, isolating and promoting the model’s \emph{intended responses and conduct} from the full range of its \textit{knowledge and competencies} is essential for developing AI systems that are reliable, effective, and manageable \citep{ouyang2022training}. Common approaches guide LMs toward human-preferred outputs via reinforcement learning (RL); in this work, we demonstrate that the objective used in RL-based methods can in fact be solved exactly using a straightforward binary cross-entropy loss, thereby streamlining the preference optimization process.

Generally, current techniques impart desirable behavioral tendencies into language models by leveraging curated collections of human preference data, which encode notions of safety and utility. This phase of learning preferences succeeds the initial, unsupervised large-scale pre-training on textual data. Although a simple strategy for aligning model behavior is through supervised fine-tuning with high-quality human-labeled responses, reinforcement learning from human or artificial feedback (RLHF/RLAIF; \citep{christiano2017deep,bai2022constitutional}) has emerged as the prevailing method. RLHF constructs a reward model by fitting to datasets of expressed human preferences, subsequently employing RL to train a model to yield outputs associated with higher rewards, while simultaneously constraining deviation from the original pre-trained distribution. Despite its effectiveness in generating sophisticated dialogue and programming models, the RLHF framework is substantially more intricate than purely supervised approaches, necessitating the training of several LMs and embedding sampling from model policies within the optimization loop—consequently resulting in significant computational demands.

In this work, we demonstrate a method for optimizing language models in alignment with human preferences without the need for explicit reward modeling or reinforcement learning approaches. We introduce \textit{{Direct Preference Optimization} (DPO)}, an algorithm that achieves the same optimization objective as traditional RLHF techniques (namely, reward maximization subject to a KL-divergence regularization), but with a design that is both easily implemented and straightforward to train. Conceptually, DPO operates by amplifying the log probability ratio of favorable responses to unfavorable ones, while simultaneously applying a per-sample adaptive weighting scheme. This mechanism avoids the model collapse that arises from relying on a naive ratio-based objective. In parallel with existing algorithms, DPO utilizes a theoretical model of preferences—such as the Bradley-Terry framework \cite{bradley1952rankanalysis}—to gauge the extent to which a reward function matches human preference data. However, in contrast to prior approaches that employ this preference model to formulate a preference loss for reward model training followed by policy training with that reward, DPO leverages a variable transformation to express the preference loss directly as a function of the policy itself. With access to a dataset comprising human-annotated preferences among model outputs, DPO directly updates the policy via a straightforward binary cross entropy loss, thereby identifying the optimal policy for a reward function implicitly fitted to the observed preference data.

The principal contribution of this study is Direct Preference Optimization (DPO), a reinforcement learning-free technique for preference-based language model training. Empirical evaluations reveal that DPO matches or surpasses existing approaches, such as RLHF methods built on PPO, across tasks including sentiment control, summarization, and conversational modeling with models containing up to 6B parameters.

\section{Related Work}

As self-supervised language models scale up, they acquire the ability to handle certain tasks with zero-shot \citep{radford2019language} or few-shot prompts \citep{gpt3,megatron,chowdhery2022palm}. Nonetheless, their effectiveness on downstream applications and their responsiveness to user instructions are often substantially boosted by additional fine-tuning on collections of instructions paired with human-authored completions \citep{mishra-etal-2022-cross,sanh2022multitask,chung2022scaling,thoppilan2022lamda}. The practice known as `instruction-tuning' allows LLMs to better generalize to novel instructions not present in the tuning dataset and generally enhances their practicality and user friendliness \citep{chung2022scaling}. Although instruction tuning has shown strong results, it is typically simpler to gather \textit{relative} assessments of answer quality from humans than to amass extensive expert-generated exemplars. Subsequent research has thus improved LLMs by fine-tuning them with datasets structured around human preference judgments, resulting in advances in fields like translation \citep{kreutzer-etal-2018-reliability}, summarization \citep{stiennon2022learning,ziegler2020finetuning}, narrative generation \citep{ziegler2020finetuning}, and compliance with instructions \citep{ouyang2022training,ramamurthy2023is}. Such methodologies begin by learning a neural network-based reward model aligned with these preference datasets—commonly applying models like the Bradley-Terry framework \citep{bradley1952rankanalysis}—then proceed to adapt the language model via reinforcement learning techniques, such as REINFORCE \citep{williams1992reinforce}, proximal policy optimization (PPO; \cite{schulman2017proximal}), or related variants \citep{ramamurthy2023is}. A related research direction exploits LLMs that have already been fine-tuned through instruction-following tasks and human feedback to automatically create further synthetic preference datasets focused on criteria like safety or non-harmfulness \citep{bai2022constitutional}, relying mainly on minimal human guidance provided in the form of textual evaluation rubrics for the model's annotations. These lines of research reflect a synthesis of two traditions: one that centers on reinforcement learning-based training of language models for diverse objectives~\citep{Ranzato2015SequenceLT,paulus2018a,wu2018learning}, and another emphasizing learning from human feedback \citep{christiano2017deep,kupcsik2018learning}. Despite the advantages of utilizing comparative human preferences, it remains a formidable practical obstacle to adapt large language models using reinforcement learning; the present work introduces a principled, theoretically grounded strategy for optimizing models against relative preference data that avoids the need for RL.

Beyond the scope of language tasks, the study of policy learning from preferences has a history in both bandit and reinforcement learning frameworks, with multiple strategies having been introduced. In contextual bandits, when preferences or action rankings substitute for reward signals, the problem is formalized as a contextual dueling bandit (CDB; \cite{yue2012karmed,dudik2015contextual}). Without access to absolute reward values, theoretical work on CDBs replaces the concept of the optimal policy with the \textit{von Neumann winner}: this is a policy guaranteed to achieve an expected win rate of at least 50\% when compared to \textit{any} competing policy \citep{dudik2015contextual}. Nevertheless, a key distinction is that preference data in CDBs is typically acquired in an online manner, whereas learning from human preferences more commonly relies on a fixed set of offline preference annotations across action pairs \citep{yan2022human}. In a related vein, \textit{preference-based RL} (PbRL) seeks to infer policy from binary preferences originating from an \textit{unknown} scoring or reward function, rather than from explicit reward feedback \citep{BusaFekete2014,ruiz2023dueling}. A range of PbRL algorithms has been proposed; although some leverage off-policy preference collections, the majority involve first learning an underlying scoring (reward) function before optimizing policy relative to that learned model \citep{jain2013learning,BusaFekete2014,christiano2017deep,sadigh2017active,kupcsik2018learning}. In contrast to these two-stage procedures, our contribution is a direct, one-stage policy optimization approach that focuses on satisfying observed preferences.

\section{Preliminaries}\label{section:prelims}

We summarize the RLHF pipeline as described in \citeauthor{ziegler2020finetuning} and subsequent work (\citep{stiennon2022learning, bai2022training, ouyang2022training}). The process generally unfolds over three key stages: 1) supervised fine-tuning (SFT), 2) collecting preferences and learning a reward model, and 3) reinforcement learning for optimization.

\textbf{SFT}: The first step of RLHF usually involves adapting a large pre-trained language model using supervised learning on curated, high-quality datasets relevant to target applications (such as conversation, summarization, and others), producing an initial model $\pi^\text{SFT}$. 

\textbf{Reward Modelling Phase}: The next step prompts the SFT-fine-tuned model with various prompts $x$, generating response pairs $(y_1, y_2)\sim \pi^\text{SFT}(y \mid x)$. Human evaluators are then asked to indicate a preference between the responses, denoted as $y_w\succ y_l \mid x$, with $y_w$ and $y_l$ identifying the favored and unfavored completions, respectively, from the sampled pair. These preferences are assumed to reflect an underlying but inaccessible reward model $r^*(y, x)$. Several models are available for representing these preferences, with the Bradley-Terry (BT) model \cite{bradley1952rankanalysis} being a frequently employed option (alternatively, if access to ranked lists is available, more comprehensive models such as the Plackett-Luce model \citep{plackett1975analysis, luce2012individual} may also be used). Within the BT framework, the probability distribution over human preferences $p^*$ is formulated as follows:

Suppose we are provided with a fixed collection of pairwise comparison data, denoted $\mathcal{D}=\bigl\{x^{(i)}, y_w^{(i)}, y_l^{(i)}\bigr\}_{i=1}^N$, drawn from the ground-truth preference distribution $p^*$. To model reward, we introduce a parameterized function $r_{\phi}(x, y)$, and the optimal parameters can be obtained via maximum likelihood estimation. Recasting the preference modeling problem as binary classification leads to the following negative log-likelihood objective:

\begin{equation}\label{eq:reward_model}
    \mathcal{L}_R(r_{\phi}, \mathcal{D}) = -\mathbb{E}_{(x, y_w, y_l)\sim \mathcal{D}}\left[\log \sigma(r_{\phi}(x, y_w)- r_{\phi}(x, y_l))\right]
\end{equation}

where $\sigma$ represents the logistic sigmoid function. For large language models (LMs), it is standard practice to initialize the reward estimator $r_{\phi}(x, y)$ with the pre-trained SFT model $\pi^\text{SFT}(y \mid x)$ and append a linear projection on top of the terminal transformer layer, yielding a scalar reward output~\cite{ziegler2020finetuning}. To reduce the variability of the learned rewards, common approaches perform reward normalization, enforcing that $\mathbb{E}_{x,y\sim \mathcal{D}}\left[r_\phi(x, y)\right] = 0$ across all $x$.

\textbf{RL Fine-Tuning Phase}: Once the reward function has been acquired, it serves as the feedback signal during reinforcement learning-based model refinement. Building on earlier work~\citep{jaques2017sequence, jaques2020human}, the objective is formulated as

\begin{equation}\label{eq:RL}
\max_{\pi_{\theta}}  \mathbb{E}_{x\sim \mathcal{D}, y\sim \pi_{\theta}(y \mid x)}\bigl[r_{\phi}(x, y)\bigr] - \beta\mathbb{D}_{\textrm{KL}}\bigl[\pi_{\theta}(y\mid x)\mid \mid \pi_\text{ref}(y\mid x)\bigr],
\end{equation}

where the hyperparameter $\beta$ modulates the extent to which the new policy $\pi_\theta$ diverges from a reference policy $\pi_\text{ref}$—typically, the SFT baseline $\pi^\text{SFT}$. Commonly, $\pi_\theta$ is also initialized using the SFT model. This KL-regularization term is crucial for confining the model’s outputs to the regime where the reward function provides valid feedback and for preserving output diversity, thereby avoiding degeneration to overly narrow high-reward sequences. Because text generation produces discrete sequences, this objective cannot be differentiated directly; consequently, reinforcement learning methods are generally used for optimization. A standard strategy~\cite{ziegler2020finetuning, stiennon2022learning, bai2022training, ouyang2022training} is to define the composite reward as $r(x, y) = r_{\phi}(x, y) -\beta (\log \pi_{\theta}(y\mid x) - \log \pi_\text{ref}(y\mid x))$ and apply PPO~\cite{schulman2017proximal} to maximize it.

\section{Direct Preference Optimization}\label{sec:DPO}

Given the significant practical and computational difficulties of using reinforcement learning in large-scale fine-tuning scenarios, we propose to develop a straightforward policy optimization method that directly exploits preference data. Distinct from standard RLHF pipelines—which separately train a reward function and then improve policy via RL—our method employs a specific reward parameterization that makes the corresponding optimal policy analytically tractable, thus removing the need for RL-driven optimization. The main contribution is to utilize an explicit analytical correspondence from the reward landscape to the optimal policy, facilitating a change of variables from the space of reward functions to the space of policies. This technique eliminates the requirement for an independently trained reward model while maintaining consistency with established models of human choice, such as the Bradley-Terry model. Effectively, this results in a policy network that simultaneously encapsulates both the language modeling and the (implicit) reward estimation.

\textbf{Deriving the DPO objective.} Our starting point is the standard RL objective, Eq.~\ref{eq:RL}, framed with a general reward function $r$, as utilized in previous studies. Building upon the methodology in~\citep{peters2007reinforcement, peng2019advantage, korbak2022reinforcement, go2023aligning}, it can be readily established that the solution which optimizes the KL-regularized reward maximization objective of Eq.~\ref{eq:RL} assumes the following form:

\begin{equation}\label{eq:op_policy}
    \pi_r(y\mid x) = \frac{1}{Z(x)}\pi_\text{ref}(y\mid x)\exp\left(\frac{1}{\beta}r(x, y)\right),
\end{equation}

where the partition function $Z(x) =\sum_{y}\pi_\text{ref}(y\mid x)\exp\left(\frac{1}{\beta}r(x, y)\right)$. Details for the complete derivation can be found in Appendix \ref{app:derivation1}. Even with the maximum likelihood estimation $r_{\phi}$ of the true reward $r^*$, evaluating the normalization constant $Z(x)$ remains computationally demanding~\citep{korbak2022reinforcement, go2023aligning}, making direct application of this parameterization challenging. Nevertheless, by manipulating Eq.~\ref{eq:op_policy}, we can reformulate the reward $r$ as a function of the optimal policy $\pi_r$, the baseline policy $\pi_\text{ref}$, and the (possibly intractable) partition function $Z(\cdot)$. In particular, taking the logarithm of both sides of Eq.~\ref{eq:op_policy} and rearranging yields:

\begin{equation}\label{eq:main_eq}
    r(x,y) =\beta \log \frac{\pi_r(y\mid x)}{\pi_\text{ref}(y\mid x)} + \beta \log Z(x).
\end{equation}

This transformation is equally applicable to the true reward $r^*$ and the associated optimal policy $\pi^*$. Conveniently, the Bradley-Terry model is determined only by reward differences between two responses; in particular, ${p^*(y_1 \succ y_2 \mid x) = \sigma(r^*(x, y_1) - r^*(x, y_2))}$. Replacing $r^*(x, y)$ in Eq.~\ref{eq:bradley-terry} with the form from Eq.~\ref{eq:main_eq} leads to a cancellation of the partition function, thus allowing the probability of human preferences to be written solely in terms of the optimal policy $\pi^*$ and reference $\pi_\text{ref}$. Consequently, under the Bradley-Terry framework, the optimal RLHF policy $\pi^*$ must satisfy:

\begin{equation}\label{eq:objective}
    p^*(y_1\succ y_2 \mid x)=\frac{1}{1 + \exp\left(\beta \log \frac{\pi^*(y_2\mid x)}{\pi_\text{ref}(y_2\mid x)} - \beta \log \frac{\pi^*(y_1\mid x)}{\pi_\text{ref}(y_1\mid x)}\right)}
\end{equation}

For full details, see Appendix~\ref{app:derivation2}. Although Eq.~\ref{eq:objective} is specific to the Bradley-Terry model, parallel derivations hold for more general Plackett-Luce models~\citep{plackett1975analysis, luce2012individual}, as presented in Appendix~\ref{app:plackett_luce_models}.

Having established the probability of observing human preference data as a function of the optimal policy rather than the reward, we can now construct a maximum likelihood estimation objective for a parameterized policy $\pi_\theta$. Mirroring the reward modeling procedure (see Eq.~\ref{eq:reward_model}), we write the objective for the policy as follows:

\begin{equation}\label{eq:optimum_model}
    \mathcal{L}_\text{DPO}(\pi_{\theta}; \pi_\text{ref}) = -\mathbb{E}_{(x, y_w, y_l)\sim \mathcal{D}}\left[\log \sigma \left(\beta \log \frac{\pi_{\theta}(y_w\mid x)}{\pi_\text{ref}(y_w\mid x)} - \

In this approach, we construct an implicit reward function through a different parameterization, with its optimal policy directly corresponding to $\pi_\theta$. Additionally, our methodology aligns with fitting a reparameterized Bradley-Terry model, endowing it with desirable theoretical characteristics such as consistency, provided that the preference data distribution satisfies certain assumptions \cite{bong2022generalized}. A more in-depth discussion of the theoretical underpinnings of DPO, including its connections to related methods, is presented in Section~\ref{sec:theory}.

\textbf{What effect does the DPO update have?} To develop an intuitive grasp of the DPO algorithm, we can examine the gradient of its loss, $\mathcal{L}_\text{DPO}$. The derivative of this loss with respect to $\theta$ takes the following form:

\begin{multline*}\label{eq:gradient}
    \nabla_\theta \mathcal{L}_\text{DPO}(\pi_\theta;\pi_\text{ref}) = \\ -\beta\mathbb{E}_{(x, y_w, y_l) \sim \mathcal{D}} \bigg[\underbrace{\sigma(\hat{r}_\theta(x, y_l) - \hat{r}_\theta (x, y_w))}_\text{larger emphasis when reward order is incorrect}\bigg[\underbrace{\nabla_\theta\log \pi(y_w \mid x)}_\text{promote $y_w$ occurrence} - \underbrace{\nabla_\theta\log\pi(y_l \mid x)}_\text{discourage $y_l$ occurrence}\bigg]\bigg],
\end{multline*}

with $\hat{r}_\theta(x, y) = \beta \log \frac{\pi_\theta(y \mid x)}{\pi_\text{ref}(y \mid x)}$ representing the implicit reward associated with the output of the language model $\pi_\theta$ and the reference model $\pi_\text{ref}$ (see further explanation in Section~\ref{sec:theory}). In essence, the gradient of $\mathcal{L}_\text{DPO}$ incentivizes an increase in the likelihood of the preferred responses $y_w$ while discouraging those of the less-preferred responses $y_l$. Crucially, this update is modulated by the extent to which the implicit reward $\hat{r}_\theta$ incorrectly ranks $y_l$ above $y_w$, scaled by the factor $\beta$; this reflects both the severity of the misordering and the influence of the KL regularization. Empirical results indicate that this weighted formulation is critical, as omitting the weighting factor—as in a simplified version of the algorithm—may result in undesirable model behaviors (see Appendix Table~\ref{tab:unlikelihood_generations}).

\textbf{DPO overview.} The standard DPO procedure involves: 1) For each prompt $x$, generate completions $y_1, y_2 \sim \pi_\text{ref}(\cdot \mid x)$, then annotate these with human preference labels to create an offline preference dataset $\mathcal{D} = \{x^{(i)}, y_w^{(i)}, y_l^{(i)}\}_{i=1}^N$; 2) Train the language model $\pi_\theta$ by minimizing $\mathcal{L}_\text{DPO}$, using the provided $\pi_\text{ref}$, the constructed dataset $\mathcal{D}$, and a chosen value of $\beta$. In applied settings, practitioners generally seek to utilize available public preference datasets, rather than carrying out new sampling and preference collection. If these datasets are generated from $\pi^\text{SFT}$, we initialize $\pi_\text{ref}$ as $\pi^\text{SFT}$. If no such $\pi^\text{SFT}$ exists, we instead define $\pi_\text{ref}$ via maximum likelihood estimation over preferred completions ${(x, y_w)}$, formally, ${\pi_\text{ref} = \argmax_{\pi}\mathbb{E}_{x, y_w \sim \mathcal{D}}\left[\log \pi(y_w \mid x)\right]}$. This approach addresses the mismatch between the inaccessible true reference distribution and the surrogate $\pi_\text{ref}$ adopted for DPO. Implementation details and information on hyperparameters appear in Appendix~\ref{app:implementation}.

\section{Theoretical Analysis of DPO} This section offers additional perspective on the DPO method, supplies theoretical justifications, and connects the strengths of DPO to weaknesses identified in actor-critic-based RLHF algorithms (notably PPO~\cite{schulman2017proximal}).

\label{sec:theory}

\subsection{Your Language Model Is Secretly a Reward Model} DPO achieves policy optimization using a single maximum likelihood loss, thus eliminating the need to explicitly fit a reward function or perform RL steps. Specifically, the objective in Eq. \ref{eq:main_eq} corresponds to a Bradley-Terry model where rewards are parameterized by $r^*(x, y) = \beta \log\frac{\pi^*_\theta(y \mid x)}{\pi_\text{ref}(y \mid x)}$, and the parameterized model $\pi_{\theta}$ is trained in a manner analogous to reward model optimization, as described in Eq. \ref{eq:reward_model}, up to a variable change. Here, we develop the theoretical basis for this parameterization, argue that it imposes no significant limitation on the space of possible reward models, and demonstrate that it facilitates precise recovery of the optimal policy. We commence by introducing an equivalence relation defined over reward functions.

\begin{definition}
Two reward functions $r(x, y)$ and $r'(x, y)$ are considered equivalent if and only if ${r(x, y)-r'(x, y) = f(x)}$ holds for some function $f$.
\end{definition}

This definition readily establishes an equivalence relation that divides the space of reward functions into distinct equivalence classes. The following lemmas formalize key consequences of this equivalence:

\begin{lemma}\label{lemma:same_prefrence} Within the Plackett-Luce and, specifically, the Bradley-Terry preference frameworks, any pair of reward functions belonging to the same equivalence class result in identical preference distributions.
\end{lemma}

\begin{lemma}\label{lemma:same_policy}
    Within the setting of constrained RL, reward functions from the same equivalence class induce the same optimal policy.
\end{lemma}

The proofs for these statements are immediate, and can be found in Appendix \ref{app:lemma1}. The first lemma highlights a classic under-identifiability challenge encountered with models in the Plackett-Luce family~\cite{plackett1975analysis}. To mitigate this ambiguity, it is standard practice to introduce additional identifiability conditions to ensure meaningful maximum likelihood estimates in Eq. \ref{eq:reward_model}~\cite{bong2022generalized}. The second lemma demonstrates that every member within an equivalence class produces the same optimal policy. Thus, for our learning objective, it suffices to recover any reward function representative from the optimal equivalence class. The following theorem, with proof in Appendix~\ref{app:thm1}, formalizes this idea:

\begin{theorem}\label{thm:main}
    Under mild conditions, every equivalence class of reward functions compatible with the Plackett-Luce (including Bradley-Terry) models is expressible via the reparameterization ${r(x, y) = \beta \log \frac{\pi(y\mid x)}{\pi_\text{ref}(y\mid x)}}$ for some model $\pi(y\mid x)$ and a fixed reference model $\pi_\text{ref}(y \mid x)$.
\end{theorem}

\begin{sproof}
    Let $r(x, y)$ denote a given reward function, associated with an optimal policy $\pi_r(y \mid x)$ as defined in Eq. \ref{eq:op_policy}. We will demonstrate that a member of the equivalence class of $r$ admits the reparameterization above. Define the transformation $f$ by
\begin{equation}
    f(r; \pi_\text{ref}, \beta)(x, y) = r(x, y) - \beta\log\sum_{y}\pi_\text{ref}(y\mid x)\exp\left(\frac{1}{\beta}r(x, y)\right)
\end{equation}
This operator $f$ re-centers the reward function by subtracting the log-partition function (which depends only on $x$). Since this adjustment involves only $x$, $f(r; \pi_\text{ref}, \beta)(x, y)$ remains in the equivalence class of $r(x, y)$. Substituting $r$ with the right-hand side of Eq.~\ref{eq:main_eq}, we see that $f(r; \pi_\text{ref}, \beta)(x, y) = \beta \log \frac{\pi_r(y\mid x)}{\pi_\text{ref}(y\mid x)}$. Thus, the projection $f$ yields a representative of $r$'s equivalence class in the desired parametric form, so our chosen reparameterization retains full generality.
\end{sproof}

An alternative interpretation of Theorem~\ref{thm:main} is that it precisely determines which reward function within each equivalence class is chosen by the DPO reparameterization—specifically, the reward function that fulfills:

\begin{equation}\label{eq:lag_p}
     \sum_{y}\underbrace{\pi_\text{ref}(y\mid x)\exp\left(\frac{1}{\beta}r(x, y)\right)}_{=\pi(y\mid x)\text{, using Thm.~\ref{thm:main} reparam.}} = 1,
\end{equation}

This condition ensures that $\pi(y\mid x)$ is a well-defined probability distribution, with non-negative probabilities summing to unity. Moreover, referring back to Eq.~\ref{eq:op_policy}, it is evident that Eq.~\ref{eq:lag_p} corresponds to the partition function for the optimal policy determined by the reward $r(x, y)$. A central insight underlying the DPO approach is the ability to enforce specific constraints on the generally under-determined Plackett-Luce (and, as a particular case, Bradley-Terry) class of preference models. This allows the set of possible reward functions to remain unchanged, yet the optimal policy described in Eq. \ref{eq:op_policy} becomes analytically manageable for every prompt $x$.

\subsection{Instability of Actor-Critic Algorithms} Our framework also facilitates the identification of instability issues in standard actor-critic methods commonly applied in RLHF, such as PPO. In line with the RLHF process, we concentrate on the RL fine-tuning stage discussed in Section \ref{section:prelims}. Parallels may be drawn with the control as inference perspective \cite{levine2018reinforcement}, especially regarding the constrained RL setup introduced in \ref{eq:RL}. Considering a parameterized policy model $\pi_{\theta}(y\mid x)$, we minimize the Kullback-Leibler divergence $\mathbb{D}_{\text{KL}}[\pi_{\theta}(y|x) \mid \mid \pi^*(y\mid x)]$, where $\pi^*$ denotes the optimal policy from Eq. \ref{eq:optimum_model} corresponding to the reward parameterization $r_{\phi}(y, x)$. Some manipulations yield the following optimization criterion:

\begin{equation}\label{eq:AC}
    \max_{\pi_{\theta}}\mathbb{E}_{\pi_{\theta}(y\mid x)}\bigg[\underbrace{r_{\phi}(x, y) -\beta\log\sum_{y}\pi_\text{ref}(y\mid x)\exp\left(\frac{1}{\beta}r_{\phi}(x, y)\right)}_{f(r_{\phi}, \pi_\text{ref}, \beta)} - \underbrace{\beta\log\frac{\pi_{\theta}(y\mid x)}{\pi_\text{ref}(y\mid x)}}_{\text{KL}}\bigg]
\end{equation}

This objective mirrors those optimized in earlier studies 
\citep{ziegler2020finetuning, stiennon2022learning, bai2022training, ouyang2022training}, employing the DPO-equivalent reward within the reward class $r_{\phi}$. In this framework, the normalization term in $f(r_{\phi}, \pi_\text{ref}, \beta)$ can be interpreted as representing the soft value function of the reference policy $\pi_\text{ref}$. Although this normalization does not impact the set of optimal solutions, omitting it can lead to a high-variance policy gradient, potentially destabilizing the learning process. It is possible to incorporate this normalization through a learned value function, but this approach is also challenging to optimize reliably. Previous research has addressed the normalization issue by using a baseline derived from human completions, effectively employing a single-sample Monte Carlo estimate of the normalization term. In contrast, the DPO reparameterization produces a reward function that eliminates the need for such baselines.

\section{Experiments}
Here, we empirically assess the effectiveness of DPO for learning policies directly from preference data. We begin by investigating a carefully controlled text generation scenario, examining how efficiently DPO navigates the tradeoff between reward maximization and KL-divergence minimization relative to standard preference learning methods such as PPO. Subsequently, we extend our evaluation to larger models and more complex RLHF benchmarks, encompassing both summarization and dialogue tasks. Our results indicate that, with minimal hyperparameter tuning, DPO frequently matches or surpasses robust baselines like RLHF with PPO and the selection of the top-performing sample among $N$ trajectories under a learned reward function. Prior to discussing these empirical findings, we outline our experimental protocols; further methodological specifics are provided in Appendix~\ref{app:exp_details}.

\textbf{Tasks.} We conduct experiments across three distinct open-ended text generation tasks. For each task, the learning algorithms derive a policy from a dataset of preferences $\mathcal{D}=\bigl\{x^{(i)}, y_w^{(i)}, y_l^{(i)}\bigr\}_{i=1}^N$. In the \textbf{controlled sentiment generation} task, $x$ denotes a movie review prefix sourced from the IMDb dataset \cite{maas-EtAl:2011:ACL-HLT2011}, with the objective for the policy being to produce a continuation $y$ exhibiting positive sentiment. To facilitate controlled measurement, we synthesize preference pairs over generated outputs using a pre-trained sentiment classifier, ensuring that $p(\text{positive}\mid x,y_w)>p(\text{positive}\mid x,y_l)$. For SFT, we conduct fine-tuning of GPT-2-large on training-set reviews from IMDb until performance stabilizes (further information provided in App~\ref{app:sentiment_details}). In the \textbf{summarization} task, $x$ corresponds to a Reddit forum post, where the model’s goal is to generate a summary $y$ highlighting the central ideas. Consistent with previous studies, we utilize the Reddit TL;DR summarization dataset \citep{volske-etal-2017-tl} and a human preference set compiled by \citeauthor{stiennon2022learning}. The SFT model is fine-tuned using human-authored forum post summaries\footnote{\url{https://huggingface.co/CarperAI/openai_summarize_tldr_sft}} and employs the TRLX \citep{leandro_von_werra_2023_7790115} library for RLHF training. The referenced human preference dataset consists of ratings on samples produced by a separate but comparably-trained SFT model, as gathered by \citeauthor{stiennon2022learning}. Lastly, in the \textbf{single-turn dialogue} scenario, $x$ represents a user’s prompt, which may range widely in subject from science to personal advice. Here, the policy must generate an informative and engaging answer $y$; we make use of the Anthropic Helpful and Harmless dialogue dataset \citep{bai2022training}, which provides 170k conversations between human users and an AI assistant. Each conversation concludes with two model-generated responses and a human-annotated preference label for the more desirable answer. As there is no existing SFT model for this dataset, we fine-tune a general-purpose language model exclusively on the preferred responses to create the SFT baseline.

\textbf{Evaluation.} Our experimental framework employs two distinct evaluation strategies. To examine how well each algorithm optimizes the constrained reward maximization goal, we analyze their performance on a frontier defined by the achieved reward and the KL-divergence relative to a reference policy within the controlled sentiment generation scenario. This is feasible because we have direct access to the authentic reward signal, provided by a sentiment classifier. In practical scenarios, where the ground truth reward is typically unavailable, we instead assess algorithm performance using their \textit{win rate} over a baseline policy. Here, GPT-4 serves as a stand-in for human assessors, evaluating the quality of summaries and the helpfulness of single-turn dialogue responses, respectively, in summarization and dialogue tasks. For the summarization task, the baseline comprises the test set’s reference summaries, whereas for dialogue, it is the response marked as preferred in the test data. Although prior work has indicated that language models can surpass traditional automatic metrics in evaluator effectiveness \citep{Chen2023ExploringTU}, we supplement our results with a human evaluation to support the use of GPT-4 as an evaluator (see Sec.~\ref{sec:human-judgments}). We observe a strong alignment between GPT-4 and human judgments, with the agreement between human annotators and GPT-4 typically on par with, or surpassing, inter-annotator reliability.

\textbf{Methods.} Alongside DPO, our comparative analysis includes a range of established methods for training language models to follow human preferences. The simplest baseline involves zero-shot prompting of \textbf{GPT-J} \citep{gpt-j} for summarization, and two-shot prompting of \textbf{Pythia-2.8B} \citep{biderman2023pythia} for dialogue. Additional baselines feature the \textbf{SFT} model and \textbf{Preferred-FT}, which applies supervised fine-tuning to the chosen completion $y_w$—sourced from either the SFT model (in the sentiment and summarization domains) or a standard language model (in dialogue tasks). We also assess the \textbf{Unlikelihood} method~\citep{welleck2019neural}, which modifies the policy by promoting the likelihood of $y_w$ and demoting the likelihood of $y_l$, governed by an optional scaling factor $\alpha\in[0,1]$ on the unlikelihood component. Furthermore, we investigate \textbf{PPO} \citep{schulman2017proximal} utilizing a reward model inferred from preferences, and \textbf{PPO-GT}, which operates with access to the ground-truth reward in the sentiment-controlled setup. For sentiment-related tasks, we implement two variants of PPO-GT: a standard version from \cite{leandro_von_werra_2023_7790115} and a modified version that introduces reward normalization and hyperparameter adjustments to boost outcomes—these refinements are also applied when using PPO with learned rewards. Lastly, we evaluate the \textbf{Best of $N$} strategy, where $N$ samples are drawn from the SFT model (or Preferred-FT for dialogue), and the candidate with the top score from a reward model trained on the preference data is selected. Despite its strong empirical results, this baseline dissociates the reward model’s capacity from PPO optimization but suffers from limited scalability due to the requirement to sample $N$ completions for every single test-time query.

\subsection{How well can DPO optimize the RLHF objective?}

Standard RLHF algorithms typically utilize a reward maximization objective with a KL constraint, ensuring the policy seeks higher rewards without straying significantly from the reference policy. Thus, a meaningful comparison of these algorithms should jointly consider both the obtained reward and the corresponding KL divergence. Attaining marginally greater rewards at the cost of a substantially increased KL is generally not favorable. Figure~\ref{fig:frontier-tldr-main} presents the reward-KL tradeoff curves for different methods in the sentiment domain. For each algorithm, we conduct several training trials, varying the hyperparameter governing policy conservativeness for each run (target KL values $\in\{3,6,9,12\}$ for PPO, $\beta \in \{0.05,0.1,1,5\}$, and $\alpha\in\{0.05,0.1,0.5,1\}$ for unlikelihood; preferred-FT is varied by random seeds). In total, the experimental sweep comprises 22 training runs. At every 100 training steps up to convergence, we assess each policy's performance on a set of test prompts, reporting both the mean reward under the true reward function and the average KL at the sequence level\footnote{This is computed as the total of per-timestep KL divergences.} with respect to the reference policy $\text{KL}\left(\pi\mid \mid \pi_\text{ref}\right)$. Our analysis reveals that DPO establishes a markedly more optimal reward-KL frontier than all other approaches, attaining superior rewards at comparable, low KL levels. This outcome is especially significant for several reasons: although both DPO and PPO seek to optimize the same objective, DPO is substantially more effective, with its reward/KL tradeoff uniformly surpassing that of PPO. Additionally, DPO produces a more advantageous frontier than PPO, \emph{even in cases where PPO has direct access to the ground truth rewards} (PPO-GT).

\subsection{Can DPO scale to real preference datasets?}
\label{sec:dpo-real-datasets}
We proceed to assess the ability of DPO to fine-tune language models using real-world preference data in the contexts of summarization and single-turn dialogue tasks. In the realm of summarization, it is well established that automatic metrics such as ROUGE may not reliably reflect human judgments~\citep{stiennon2022learning}. Previous research has demonstrated the effectiveness of PPO-based fine-tuning on human preference data for generating superior summaries. In our experiments, we generate samples from the test partition of the TL;DR summarization dataset and calculate the mean win rate when comparing these completions to reference outputs from the test set. Each method generates completions across a range of sampling temperatures between 0.0 and 1.0; the resulting win rates are visualized in Figure~\ref{fig:frontier-tldr-main} (right). For consistency, DPO, PPO, and Preferred-FT all build upon the same GPT-J SFT model\footnote{\url{https://huggingface.co/CarperAI/openai_summarize_tldr_sft}} during fine-tuning. Our findings indicate that DPO attains a win rate near 61\% at a temperature of 0.0, outperforming PPO, which achieves approximately 57\% at its own optimal temperature setting. Furthermore, DPO's maximum win rate surpasses that achieved by the $N$-best baseline method. It is important to mention that DPO's $\beta$ hyperparameter was not subject to extensive tuning, suggesting these results may not fully capture DPO's capabilities. Additionally, we observe that DPO exhibits significantly greater resilience to changes in sampling temperature, whereas PPO's win rate may drop to that of the underlying GPT-J model at higher temperatures. In contrast, Preferred-FT provides minimal improvement over the original SFT model. Finally, a direct human comparison between DPO and PPO, detailed in Section~\ref{sec:human-judgments}, reveals that DPO samples (temperature 0.25) are chosen over PPO samples (temperature 0) in 58\% of cases.

For single-turn dialogue evaluation, we examine the performance of the various approaches using the subset of the Anthropic HH dataset \citep{bai2022training} test split that includes only one exchange between human and assistant. In assessments involving GPT-4, we measure the win rate for each method based on how often they produce completions that match the preferred responses in the test set. Due to the absence of an established SFT model for this scenario, we begin with a pre-trained Pythia-2.8B and train a reference model via Preferred-FT on selected completions to ensure output distribution alignment. This reference model is subsequently fine-tuned with DPO. For baselines, we compare against the top-performing completion among 128 samples from Preferred-FT (noting that improvements for the Best of $N$ baseline saturate at $N=128$ for this dataset, as demonstrated in Appendix Figure~\ref{fig:best-of-n}) and against the Pythia-2.8B base model configured with a 2-shot prompt; in these experiments, DPO matches or outperforms the strongest results at optimal temperature settings for each technique. Furthermore, we consider an RLHF-trained model using PPO on the Anthropic HH dataset \footnote{\url{https://huggingface.co/reciprocate/ppo_hh_pythia-6B}} obtained from a widely recognized repository \footnote{\url{https://github.com/CarperAI/trlx/tree/main/examples/hh}}, but cannot identify any prompt or temperature configuration that surpasses the base Pythia-2.8B’s performance. Given our findings from TL;DR and recognizing both Best of 128 and PPO optimize the identical reward function, we regard Best of 128 as an approximate indicator of PPO performance. In summary, DPO stands out as the sole method that is both computationally efficient and capable of surpassing the preferred completions within the Anthropic HH dataset, while achieving comparable or superior outcomes relative to the more resource-intensive Best of 128 approach. Additionally, Figure~\ref{fig:dialogue-main} highlights the rapid convergence of DPO to its peak performance.

\subsection{Generalization to a new input distribution}

To more thoroughly investigate the robustness of PPO and DPO under distributional shifts, we assess the models trained in the Reddit TL;DR summarization setup by testing their policies on news articles from the test split of the CNN/DailyMail dataset \citep{nallapati-etal-2016-abstractive}. For consistency, we adopt the best sampling temperatures from TL;DR (0 and 0.25). Table~\ref{tab:ood} summarizes the outcomes. Using an adapted version of the GPT-4 (C) prompt—modified to substitute “news article” in place of “forum post”—we measure the GPT-4 win rate compared to the datasets’ reference summaries. DPO maintains a substantial advantage over PPO when confronted with this distributional change. These findings offer preliminary support that DPO methods can generalize comparably to PPO approaches, even though DPO lacks access to the additional, unlabeled Reddit TL;DR prompts utilized by PPO.

\subsection{Validating GPT-4 judgments with human judgments}
\label{sec:human-judgments}
To establish the credibility of GPT-4’s preference assessments, we carry out a human evaluation based on the TL;DR summarization task, employing two distinct GPT-4 prompting strategies. The \textbf{GPT-4 (S)} (“simple”) prompt directs GPT-4 to judge which summary better conveys key information from the post. In contrast, the \textbf{GPT-4 (C)} (“concise”) prompt adds an emphasis on conciseness, which we include after observing that the “simple” prompt leads GPT-4 to favor summaries that are longer and more repetitive than those favored by human evaluators. The full text of the prompts can be found in Appendix~\ref{app:prompts}. Our evaluation spans three conditions—top-performing (DPO, temp. 0.25), lowest-performing (PPO, temp. 1.0), and an intermediate (SFT, temp. 0.25) approach—each compared against the PPO model using its optimal greedy sampling configuration. This design enables us to sample a broad spectrum of generation qualities; notably, each is contrasted with the PPO baseline. Our analysis reveals that across both prompting variants, GPT-4's preferences align with those of human raters at approximately the same rate as human-to-human agreement, indicating that GPT-4 serves as a reasonable proxy for manual evaluations (note that, owing to a limited pool of human annotators, only the DPO and PPO-1 matchups feature repeated human assessments). Generally, we observe that win rates with the \textbf{GPT-4 (C)} prompt are more consistent with human choices, and thus adopt it for reporting main results in Section~\ref{sec:dpo-real-datasets}. Further information about the human evaluation protocol—including the rater interface and the list of volunteer raters—can be found in Appendix~\ref{app:human-study}.

\section{Discussion}
Preference-based learning presents a robust and adaptable approach for developing competent and value-aligned language models. In this work, we introduced DPO, a straightforward methodology for language model training from preference data, eschewing the need for reinforcement learning. Unlike traditional methods that adapt preference learning into a conventional RL context to leverage standard RL tools, DPO establishes a direct correspondence between policy parameters in language models and associated reward functions. This allows us to train models to adhere to human preferences using only a basic cross-entropy loss, thereby removing the dependence on reinforcement learning techniques while preserving generality. Remarkably, DPO achieves performance on par with or superior to prevailing RLHF strategies, such as those utilizing PPO, and it does so with minimal hyperparameter adjustment. As a result, DPO substantially lowers the obstacles to training language models on human preferences.

\textbf{Limitations \& Future Work.} The findings in this paper open several avenues for subsequent investigation. One open question concerns the extent to which the DPO-trained policy is able to generalize to distributions outside the training data, especially in comparison with explicit reward-based learning. Preliminary evidence indicates that DPO achieves a level of generalization comparable to models trained with PPO, yet this warrants more exhaustive examination. Another area for exploration involves determining whether DPO-based self-labeling can be effectively leveraged with prompts lacking human annotations. Furthermore, the phenomenon of reward over-optimization in the context of direct preference optimization—potentially reflected in the modest reduction in performance seen in Figure~\ref{fig:dialogue-main}-right—remains to be fully characterized. Additionally, although our evaluation includes models as large as 6B parameters, there is significant promise in extending DPO to much larger, state-of-the-art architectures. On the topic of assessment, we observed that GPT-4-generated win rates are sensitive to prompting strategies; future research might identify more reliable methods for obtaining robust automated model judgments. Lastly, DPO's framework lends itself to a variety of tasks beyond language modeling with human preference data, including the training of generative models in diverse domains.